\documentclass[../BTL.tex]{subfiles}
\begin{document}
\chapter{Tham khảo cấu hình bài kiểm thử}
\label{phu-luc:chapter}

Phụ lục này tóm tắt các loại bước kiểm thử, assertion và cách khai báo biến trích xuất; mỗi loại có ví dụ ngắn để tra cứu khi tạo kế hoạch chấm.

\section{Cấu trúc kế hoạch và bước chấm}

Một kế hoạch gồm các bước chạy theo thứ tự \textnormal{sequenceOrder}. Mỗi bước có loại, cấu hình, trọng số và tùy chọn dừng kế hoạch khi thất bại.

\begin{table}[H]
\centering
\begin{tabularx}{\textwidth}{|p{3.1cm}|X|}
\hline
Trường & Ý nghĩa / ví dụ \\ \hline
\textnormal{sequenceOrder} & Thứ tự chạy kế hoạch, ví dụ \textnormal{1}. \\ \hline
\textnormal{weight} & Trọng số kế hoạch hoặc bước, ví dụ \textnormal{1}. \\ \hline
\textnormal{stepType} & Loại bước: \textnormal{HTTP\_REQUEST}, \textnormal{DB\_QUERY}, \textnormal{DB\_SCHEMA\_CHECK}, hoặc \textnormal{DB\_MIGRATION}. \\ \hline
\textnormal{config} & Cấu hình riêng theo loại bước. \\ \hline
\textnormal{required} & \textnormal{true} để dừng các bước còn lại trong kế hoạch khi bước này thất bại. \\ \hline
\end{tabularx}
\caption{Các trường cơ bản của kế hoạch và bước chấm.}
\label{tab:plan-step-quick-reference}
\end{table}

\section{Bước HTTP\_REQUEST}

Cấu hình thường dùng gồm phương thức, đường dẫn, dữ liệu gửi đi, trạng thái mong đợi, assertions và extract. Có thể thêm \textnormal{headers}, \textnormal{query\_params}, \textnormal{data} hoặc \textnormal{timeoutMs} khi cần.

\begin{table}[H]
\centering
\begin{tabularx}{\textwidth}{|p{3.1cm}|X|}
\hline
Trường & Ví dụ / ý nghĩa \\ \hline
\textnormal{method} & \textnormal{GET}, \textnormal{POST}, \textnormal{PUT}, \textnormal{PATCH}, hoặc \textnormal{DELETE}. \\ \hline
\textnormal{path} & Đường dẫn API, ví dụ \textnormal{/api/v1/books}. \\ \hline
\textnormal{body} & JSON gửi trong yêu cầu, thường dùng với \textnormal{POST} hoặc \textnormal{PUT}. \\ \hline
\textnormal{expected\_status} & Mã HTTP mong đợi, ví dụ \textnormal{200} hoặc \textnormal{201}; có thể thay bằng assertion \textnormal{status}. \\ \hline
\textnormal{assertions} & Danh sách điều kiện kiểm tra phản hồi; mặc định mọi điều kiện đều phải đúng. \\ \hline
\textnormal{extract} & Danh sách giá trị lấy từ phản hồi để dùng ở bước sau. \\ \hline
\end{tabularx}
\caption{Các trường thường dùng trong bước HTTP.}
\label{tab:http-step-quick-reference}
\end{table}

\subsection{Các loại assertion}

Tên \textnormal{kind} phân biệt chữ hoa và chữ thường. Các assertion trong cùng danh sách dùng phép AND; \textnormal{any\_of} cho phép gom các điều kiện theo phép OR. Ví dụ là các trường còn lại của object assertion; giá trị \textnormal{kind} lấy theo cột Loại.

\begin{table}[H]
\centering
\small
\begin{tabularx}{\textwidth}{|p{3.5cm}|X|>{\raggedright\arraybackslash}p{4.3cm}|}
\hline
Loại & Dùng để kiểm tra & Ví dụ \\ \hline
\textnormal{status} & Mã trạng thái HTTP & \texttt{equals: 201} \\ \hline
\textnormal{contains} & Body có chứa chuỗi & \texttt{text: "created"} \\ \hline
\textnormal{json\_path} & Đường dẫn JSON tồn tại hoặc không tồn tại & \texttt{path: "\$.id", exists: true} \\ \hline
\textnormal{body\_equals} & Toàn bộ body bằng JSON kỳ vọng & \texttt{json: \{"id":"b1"\}} \\ \hline
\textnormal{body\_structure} & Tên trường và kiểu dữ liệu giống JSON mẫu & \texttt{json: \{"id":""\}} \\ \hline
\textnormal{field\_equals} & Giá trị tại đường dẫn bằng giá trị kỳ vọng; hỗ trợ biến & \texttt{path: "\$.id", equals: "\$\{id\}"} \\ \hline
\textnormal{array\_length} & Độ dài mảng tại đường dẫn & \texttt{path: "\$", equals: 2} \\ \hline
\textnormal{contains\_item} & Kết quả tại đường dẫn có phần tử cần tìm & \texttt{path: "\$[*]", equals: "A"} \\ \hline
\textnormal{json\_path\_equals} & Giá trị tại đường dẫn bằng giá trị kỳ vọng; hỗ trợ biến & \texttt{path: "\$.id", equals: "\$\{id\}"} \\ \hline
\textnormal{greater\_than} & Giá trị số lớn hơn ngưỡng & \texttt{path: "\$.price", value: 100} \\ \hline
\textnormal{less\_than} & Giá trị số nhỏ hơn ngưỡng & \texttt{path: "\$.price", value: 500} \\ \hline
\textnormal{matches\_regex} & Giá trị khớp biểu thức chính quy & \texttt{path: "\$.id", pattern: "ABC[0-9]+"} \\ \hline
\textnormal{any\_of} & Ít nhất một assertion con đúng (OR) & \texttt{assertions: [contains(A), contains(B)]} \\ \hline
\end{tabularx}
\caption{Các loại assertion được hỗ trợ và ví dụ.}
\label{tab:assertion-kinds}
\end{table}

\section{Trích xuất và dùng biến}

Khai báo \textnormal{extract} để lấy giá trị từ body phản hồi HTTP. \textnormal{expression} là JsonPath; \textnormal{fallback} là giá trị thay thế nếu không tìm thấy dữ liệu.

\begin{lstlisting}[caption={Ví dụ trích xuất mã sách và dùng ở bước sau},label={lst:extract-example}]
"extract": [
  {
    "name": "bookId",
    "from": "response_body",
    "expression": "$.id"
  }
]

"path": "/api/v1/books/${bookId}"
\end{lstlisting}

Các dạng JsonPath thường dùng: \textnormal{\$.id} lấy trường ở gốc, \textnormal{\$.author.name} lấy trường lồng nhau, \textnormal{\$.items[0]} lấy phần tử đầu mảng, và \textnormal{\$.items[*]} lấy các phần tử trong mảng. Dùng \textnormal{as: "json"} để lưu đối tượng hoặc mảng thay vì một giá trị đơn.

\section{Các bước cơ sở dữ liệu}

Ba loại bước DB dùng chung thông tin kết nối. \textnormal{database}, \textnormal{username} và \textnormal{password} là bắt buộc; \textnormal{db\_type} mặc định là \textnormal{postgres}.

\begin{table}[H]
\centering
\begin{tabularx}{\textwidth}{|p{3.2cm}|X|}
\hline
Loại bước & Cấu hình chính \\ \hline
\textnormal{DB\_QUERY} & \textnormal{query}; có thể kiểm tra \textnormal{expected.row\_count} và \textnormal{expected.columns}. \\ \hline
\textnormal{DB\_\allowbreak SCHEMA\_\allowbreak CHECK} & Danh sách \textnormal{checks}: \textnormal{TABLE\_EXISTS}, \textnormal{COLUMN\_EXISTS}, \textnormal{INDEX\_EXISTS}, hoặc \textnormal{PRIMARY\_KEY}. \\ \hline
\textnormal{DB\_\allowbreak MIGRATION} & Danh sách lệnh SQL trong \textnormal{statements}. \\ \hline
\end{tabularx}
\caption{Các loại bước kiểm thử cơ sở dữ liệu.}
\label{tab:db-step-types}
\end{table}

\begin{lstlisting}[caption={Ví dụ kiểm tra bảng và khóa chính},label={lst:db-schema-example}]
{
  "connection": {
    "db_type": "postgres",
    "database": "bookstore",
    "username": "grader",
    "password": "secret"
  },
  "checks": [
    {"kind": "TABLE_EXISTS", "table_name": "books"},
    {"kind": "PRIMARY_KEY", "table_name": "books", "column": "id"}
  ]
}
\end{lstlisting}

\end{document}
